% ============================================
% DOC-07 (EN)
% English translation — SPEC-935-R355
% ============================================
\section*{\textbf{DOC-07} --- Psychiatric Diagnosis (1979)}
\markboth{DOC-07 --- Psychiatric Diagnosis (1979)}{}

\noindent \textit{DSM-II --- Dr. Heitor Oliveira}

\noindent \textbf{Patient:} Joaquim Antônio Correia --- Record 1,260
\textbf{Physician:} Dr. Heitor Oliveira, CRM-MG 4,892
\textbf{Date:} April 14, 1979

\vspace{0.5em}\noindent\rule{\textwidth}{0.5pt}\vspace{0.5em}

\noindent \textbf{ICD-8 / DSM-II}

\noindent \textbf{Axis I --- Principal Diagnosis:}
- \textbf{295.3 --- Schizophrenia, Paranoid Type}
  - Justification: Systematized persecutory delusions (believes a \textquotedblleft{}creature\textquotedblright{} inhabits him), persistent visual and tactile hallucinations, disconnected speech with neologisms (\textquotedblleft{}molambudo\textquotedblright{} as an existential state), affective blunting interspersed with outbursts of aggressiveness.

\noindent \textbf{Axis II --- Secondary Diagnosis:}
- \textbf{301.7 --- Psychopathic Personality}
  - Justification: Inability to feel remorse, recurrent aggressive conduct (three confirmed homicides in his youth), manipulation of health professionals, absence of affective bonds.

\noindent \textbf{Axis III --- Associated Medical Conditions:}
- Chronic jaundice (right eye --- possible sequela of severe childhood malnutrition, or untreated liver disease).
- Protein-calorie malnutrition consistent with the fictional chronology: the 1915 drought and the 1917 transfer.

\noindent \textbf{Axis IV --- Severity of Psychosocial Stressors:}
- Extreme.

\noindent \textbf{Axis V --- Highest Level of Adaptive Functioning in the Past Year:}
- 10/100 --- Persistent danger to self and others.

\vspace{0.5em}\noindent\rule{\textwidth}{0.5pt}\vspace{0.5em}

\noindent \textbf{Clinical Observations:}

\noindent The patient was admitted to the Hospital-Colony in 1917, at age 10. Historical records indicate that he arrived \textquotedblleft{}in a state of starvation, from the concentration camps of Ceará, orphaned of father and mother, with a yellow eye and not speaking a single word other than 'fooood'.\textquotedblright{}

\noindent Over 62 years of institutionalization, the patient presented:
- Recurring episodes of self-harm (he scratched his own face until it bled);
- Three confirmed homicides (1922, 1937, 1948) --- all victims were staff members or patients who tried to get close;
- Periods of selective mutism (he refused to speak to staff; he once went 7 years without addressing employees --- from 1949 to 1956);
- Compulsive writing of a \textquotedblleft{}diary\textquotedblright{} that he never allowed any doctor to read in full;
- Fixation on the theme of \textquotedblleft{}hunger\textquotedblright{} and of a \textquotedblleft{}creature with yellow eyes\textquotedblright{} that accompanies him.

\noindent The patient demonstrates preserved intelligence --- possibly above average --- but entirely diverted toward the construction of his delusion. His ability to narrate his own story with internal coherence (within the delusional system) is remarkable.

\noindent \textbf{Prognosis:} Guarded. Patient unrecoverable. Keep in isolation with continuous sedation.

\vspace{0.5em}\noindent\rule{\textwidth}{0.5pt}\vspace{0.5em}

\noindent \textit{Dr. Heitor Oliveira}
\textit{Attending Psychiatrist --- Maximum Security Ward}
\textit{Barbacena Hospital-Colony}

\vspace{0.5em}
\noindent \textit{P.S.: Dr Heitor Oliveira disappeared on 13 June 1979, two months after signing this report. The missing-person case was never closed. In the hospital's personnel file, a colleague noted, by hand, in the margin: \textquotedblleft{}He stopped eating three weeks earlier. He said he no longer felt hunger --- and that, he said, was the worst symptom of all.\textquotedblright{}}

\noindent\textit{\textbf{↪ Routes:} \rota{DOC-08} \rotasep \rota{DOC-04}}
